\section{Resultados históricos e contexto M5}
\label{app:historical}

Este apêndice preserva a reanálise das matrizes anteriores. Seu objetivo é documentar os antecedentes da avaliação temporal, não acrescentar observações ao teste principal. As mesmas 145 séries da Bahia e 2.408 séries de alimentos do M5 foram avaliadas com doze políticas. Características, preparação, episódios e orçamentos diferem do protocolo principal. Assim, a diferença entre custos históricos e atuais não estima o efeito isolado de uma mudança de método.

O M5 histórico agrega vendas em 93 blocos de 21 dias, incluindo um bloco final parcial. O filtro local FOODS com CV de 1,5 por produto, sobre lojas e ciclos, e a suavização superior por intervalo interquartil usam a trajetória completa. O recorte foi inspirado em Zabraoui et al.\ \cite{sca_27e994c54848}, que não publicam identificadores e corte numérico suficientes para certificar uma amostra idêntica. A nova execução M5 utiliza protocolo próprio e permanece em andamento.

\begin{table}[htbp]
\centering\small
\caption{Coortes verificadas nas matrizes analisadas. Prefixo indica o corte de ajuste configurado no código, não uma garantia de independência temporal de toda a preparação.}\label{tab:cohort}
\begin{tabular}{lrr}
\toprule
Característica & Rede brasileira (BA) & M5 FOODS\\
\midrule
Séries & 145 & 2.408\\
Lojas & 115 & 10\\
Produtos & 9 & 287\\
Ciclos disponíveis & 38 & 93\\
Prefixo configurado & 17 & 19\\
Prefixo sem demanda & 0 & 2.213\\
Menos de 3 ciclos positivos no prefixo & 0 & 2.237\\
Sparse High Impact & 116 & 2.350\\
Unstable Trend & 18 & 40\\
High Volatility Seasonal & 11 & 18\\
Políticas na matriz & 12 & 12\\
\bottomrule
\end{tabular}
\end{table}

Os prefixos configurados têm 17 ciclos na Bahia e 19 no M5. No M5, 2.213 séries (91,90\%) não registram demanda no prefixo e todas têm demanda posteriormente. Isso não comprova lançamento dos produtos ou indisponibilidade real de histórico. As diferenças de inicialização, orçamento e frequência limitam a comparação entre bases. Na dinâmica histórica, $L=2$ corresponde a chegada em $t+3$; na avaliação principal, o prazo efetivo é de dois ciclos. A heurística EOQ histórica também usa uma escala de lote distinta da EOQ por ciclo.

As matrizes registram cinco avaliações nominais por série e política, com desvios de custo praticamente nulos. Essas avaliações não são cinco treinamentos independentes. O portfólio completo da Tabela~\ref{tab:portfolio_dual} mostra que baixo custo pode acompanhar atendimento insuficiente; seus denominadores relativos são específicos de cada base.

\begin{table}[htbp]\centering\small
\caption{Portfólio histórico completo: custo relativo à melhor global viável em $\alpha=0{,}70$ e serviço médio. BA: denominador VendorResponsive; M5: denominador SA.}\label{tab:portfolio_dual}
\begin{tabular}{@{}lrrrr@{}}\toprule
 & \multicolumn{2}{c}{BA} & \multicolumn{2}{c}{M5}\\
Política & Custo relativo & Serviço & Custo relativo & Serviço\\\midrule
BigDataNewsvendor & 1,343 & 0,717 & 0,322 & 0,096\\
CappedBaseStock & 1,297 & 0,642 & 0,297 & 0,098\\
DQN & 0,292 & 0,225 & 0,318 & 0,076\\
EOQ & 1,450 & 0,942 & 0,320 & 0,133\\
FixedInterval & 1,057 & 0,844 & 0,301 & 0,100\\
GA & 2,920 & 0,896 & 0,381 & 0,231\\
GA--DQN & 2,714 & 0,731 & 0,406 & 0,467\\
Newsvendor & 0,669 & 0,540 & 0,285 & 0,084\\
PPO & 11,799 & 0,963 & 1,139 & 0,985\\
SA & 2,649 & 0,865 & 1,000 & 0,852\\
VendorResponsive & 1,000 & 0,774 & 0,287 & 0,100\\
$(s,S)$ & 1,511 & 0,943 & 0,321 & 0,134\\
\bottomrule\end{tabular}\end{table}

\subsection{Seleção retrospectiva entre lojas}
Na reanálise, a política global e as escolhas por perfil são determinadas excluindo a loja avaliada, mas consultando as saídas históricas já armazenadas. A exclusão da loja não elimina informação futura incorporada aos perfis ou à geração das matrizes. Esse desenho difere da escolha na validação e avaliação em teste posterior empregadas no corpo principal.

\begin{table}[htbp]
\centering\small
\caption{Seleção com exclusão da loja avaliada, em $\alpha=0{,}70$. A coluna abaixo do limiar conta séries individualmente, não médias de perfil. Custos não são comparáveis entre bases.}\label{tab:strategies}
\begin{tabular}{llrrr}
\toprule
Base & Estratégia & Custo médio & Fill rate & Abaixo do limiar\\
\midrule
Brasileira & Global & 433,31 & 0,7738 & 40\\
Brasileira & Perfil & 435,80 & 0,7816 & 39\\
M5 & Global & 19.500,54 & 0,8521 & 270\\
M5 & Perfil & 19.367,73 & 0,8509 & 274\\
\bottomrule
\end{tabular}
\end{table}

Em 0,70, os perfis aumentam o custo em 0,576\% na Bahia e reduzem em 0,681\% no M5, com diferenças de atendimento de 0,779 e $-0{,}122$ ponto percentual. Os intervalos condicionais de redução de custo são $[-1{,}442;0{,}153]\%$ e $[-1{,}190;2{,}402]\%$, respectivamente. Usam 5.000 reamostragens pareadas de lojas, sem reajustar escolhas; no M5, existem apenas dez lojas. Não demonstram equivalência nem vantagem geral dos perfis. A Tabela~\ref{tab:strategies} informa também as violações individuais, que não desaparecem quando a média do grupo é viável.

A sensibilidade histórica das Figuras~\ref{fig:dualmain} e~\ref{fig:dualsensitivity} ilustra mudanças de vencedor com o serviço exigido. No M5, ambas as estratégias convergem para PPO em 0,90 e 0,95. Na Bahia, o aparente ganho de custo em 0,95 acompanha menor atendimento e escolhas de contingência em grupos sem alternativa viável. Esses diagnósticos motivam a separação temporal adotada no estudo principal; não determinam um limiar comercial ótimo.

\begin{figure}[htbp]
\centering
\includegraphics[width=\textwidth]{Figure_dual_1.png}
\caption{Custo e atendimento das políticas e estratégias nas duas bases. Os painéis possuem escalas próprias: diferenças absolutas de custo entre Brasil e M5 não têm interpretação econômica comum.}
\label{fig:dualmain}
\end{figure}
\begin{figure}[htbp]
\centering
\includegraphics[width=\textwidth]{Figure_dual_2.png}
\caption{Contraste de seleção por perfil e global com exclusão de loja em diferentes limiares. Reduções de custo são interpretadas com diferenças de atendimento; intervalos são condicionais ao procedimento e às matrizes.}
\label{fig:dualsensitivity}
\end{figure}
\FloatBarrier
